\section{问题重述}

\subsection{问题背景}

反作用飞轮是卫星姿态控制系统中的关键执行机构，其轴承润滑状态直接影响长期在轨可靠性。随着润滑剂挥发、分解和磨损累积，轴承摩擦力矩逐渐增大，电机维持额定转速所需电流随之上升。因此，飞轮电流可作为反映轴承退化的可观测健康指纹。题目同时提供了飞轮仿真全寿命数据、在轨监测数据以及 XJTU-SY 轴承振动全寿命数据，为建立兼具物理解释、数据驱动特征提取和跨域迁移能力的健康管理模型提供了条件。

\subsection{问题描述}

针对题面要求，本文研究以下三个问题：

\begin{enumerate}
    \item \textbf{问题一}：根据飞轮电流、理论电流、温度和转速等历史观测，建立飞轮退化模型，判断附件 2 当前健康阶段并预测剩余寿命。
    \item \textbf{问题二}：针对 XJTU-SY 轴承振动数据，提取能够表征退化过程的特征，构造健康指数并建立轴承退化模型。
    \item \textbf{问题三}：将轴承退化规律迁移到飞轮健康预测中，修正附件 2 的 RUL 估计，并据此建立多级预警与健康管理规则。
\end{enumerate}
